\documentclass[11pt]{article}
\usepackage[utf8]{inputenc}
\usepackage[T1]{fontenc}
\usepackage{lmodern}
\usepackage{textcomp}
\usepackage{amsmath, amssymb, amsthm}
\usepackage{graphicx}
\usepackage{listings}
\usepackage[hidelinks]{hyperref}

\theoremstyle{plain}
\newtheorem{theorem}{Teorema}[section]
\newtheorem{lemma}[theorem]{Lema}
\newtheorem{proposition}[theorem]{Proposición}
\newtheorem{corollary}[theorem]{Corolario}

\theoremstyle{definition}
\newtheorem{definition}[theorem]{Definición}


\title{c7}
\date{}
\begin{document}
\maketitle

\section{Chapter 7}

\section{Forms of Argument}

A mathematical theory begins with a collection of axioms or postulates. These are statements that are assumed to be true---no verification or justification is required. The axioms are the building blocks of a theory: from them we deduce other true statements called theorems, and a proof is the process of deduction that establishes a theorem. As more and more theorems are proved, the theory is enriched by a growing list of true statements.

When we develop a mathematical argument, we need to organise sentences so that every sentence is either an axiom or a true statement derived from axioms or earlier statements. In practice only significant statements deserve to be called theorems. In the proof of a theorem many true statements may be derived, but these fragments are not usually assigned any formal label. Sometimes a proof rests on statements which, if substantial, are not of independent significance: they are called lemmas. Finally, a proposition is a statement which deserves attention, but which is not sufficiently general or significant to be called a theorem.

In this chapter we survey some methods to give shape to a mathematical argument. Our survey will continue in Chap. 8, where we consider induction techniques.

\section{7.1 Anatomy of a Proof}

A first analysis course begins with a list of axioms defining the real number system. Then one may be asked to prove that

$\forall$x $\in$R, $-$($-$x) = x. (7.1)

But isn't this statement obvious? In recognising the truth of this assertion we make implicit use of knowledge derived from experience, which is not part of the axioms. A proof from the axioms requires erasing all previous knowledge.

© Springer-Verlag London 2014 F. Vivaldi, Mathematical Writing, Springer Undergraduate Mathematics Series, DOI 10.1007/978-1-4471-6527-9\_7

113

114 7 Forms of Argument

We now introduce some axioms, state a theorem equivalent to (7.1), and then prove it from the axioms. We put the proof under X-ray. Our purpose is to dissect a mathematical argument articulating every step; we use the language mechanically, to facilitate the identification of the logical elements of the proof. To avoid making implicit assumptions, we represent familiar objects with an unfamiliar notation.

We are given a set $\Omega$ and a binary operator `' on $\Omega$, with the following properties:

G1: $\forall$x, y, $\in$$\Omega$, x y $\in$$\Omega$ G2: $\forall$x, y, z $\in$$\Omega$, (x y) z = x (y z) G3: $\exists$$\in$$\Omega$, $\forall$x $\in$$\Omega$, x = x = x G4: $\forall$x $\in$$\Omega$, $\exists$x$'$ $\in$$\Omega$, x x$'$ = 

(These axioms define a general object: a group.) Setting $\Omega$ = R, and `' = `+', one recognises that represents 0 and x$'$ represents $-$x. We are ready to state and prove our theorem.

1. Theorem.

2. $\forall$x $\in$$\Omega$, x$'$ $'$ = x.

3. Proof.

4. Let x $\in$$\Omega$ be given;

5. then (by G4) x$'$ $'$ $\in$$\Omega$;

6. let a := (x x$'$) x$'$ $'$, b := x (x$'$x$'$ $'$);

7. hence (by G1 and 5) a, b, $\in$$\Omega$, and (by G2) a = b;

8. then (by G3 and G4) a = x$'$ $'$ = x$'$ $'$;

9. then (by G3 and G4) b = x = x;

10. hence (by 7, 8, 9) x$'$ $'$ = x.

11. 

Items 1, 3, 11 are logical tags---expressions or symbols that say things about the text. They announce the statement of the theorem, and the beginning and the end of the proof, respectively. The symbol `' in item 11 may be replaced by the acronym Q.E.D.1

Items 5, 8, 9 begin with the logical tag `then', used here to mean that what comes after is deduced from the axioms. Formally, item 5 is our first theorem, being a true statement deduced from the axioms. Items 7 and 10 are deductions of a slightly different nature, and we have flagged them with the tag `hence'. These items use not only axioms but also facts assembled from previous statements.

Items 4 and 6 are not statements but instructions. Item 4 instructs us to assume that x is an arbitrary element of a set. This is a standard opening sentence which mirrors the expression $\forall$x $\in$$\Omega$ in the theorem. Item 6 instructs us to define two new quantities. Many different instructions may be found in proofs, such as the instruction to draw a picture.

1 Abbreviation for the Latin quod erat demonstrandum, meaning `which had to be demonstrated'.

7.1 Anatomy of a Proof 115

In item 10 we use implicitly the transitivity of the equality operator (if x = y and y = z, then x = z), which depends on the fact that equality is an equivalence relation (see Sect. 4.6).

$\surd$

2--- is quite removed from the axioms of arithmetic, as it relies on several definitions and facts without justification. In this respect this proof is more typical than the previous one, even though the style and conventions are the same.

Proofsfromaxiomsarenecessarilyrare.Ournextproof---theirrationalityof

1. Theorem. $\surd$ 2.

2 is irrational.

3. Proof. $\surd$

2 is rational; $\surd$

4. Suppose

5. then $\exists$m, n $\in$Z such that

2 = m/n.

6. Choose m, n in 5 so that they are co-prime. $\surd$

2 = m/n,

7. Since

8. then 2 = m2/n2;

9. then 2n2 = m2;

10. then m2 is even;

11. hence (by 10 and Theorem A2) m is even;

12. then m = 2h for some integer h;

13. hence (by 9 and 12) 2n2 = 4h2;

14. then n2 = 2h2;

15. then n2 is even;

16. hence (by 15 and theorem A) n is even;

17. hence (by 11 and 16) m and n are not co-prime;

18. hence (by 6 and 17) assumption 4 is false;

19. then its negation is true.

20. 

A fair amount of mathematics is assumed in the proof:

$\surd$

the definition of the rational numbers and of

2 (items 2, 4, 5); the definition of co-primality (item 6); some properties of equations (items 8, 9); the definition of an even integer (items 12, 17); a theorem of arithmetic, referred to as Theorem A (items 11, 16).

Theargumentcontainsnovelelements.Thecoreoftheproofisitem4,theassumption that

$\surd$ 2 is rational. This is puzzling: how can we assume that

$\surd$

2 is rational when we know it isn't? How can we let

$\surd$

2 = m/n when we haven't been told what numbers m and n are? This is not so strange; assumptions belong very much to common reasoning.

Suppose that I meet Einstein on top of Mount Everest. Suppose that people walk laterally.

2 If a prime divides the product of two integers, then it divides one of the factors.

116 7 Forms of Argument

We are clearly free to explore the logical consequences of these assumptions. In mathematics, assumptions are handled formally via the implication operator $\Rightarrow$developed in Chap. 4.

$\surd$

2isrationaleventuallyleadstoa contradiction:items6and 17 are conflicting statements. From this fact we deduce that the assumption in item $\surd$ 4 is false (item 18). Here we abandon the assumption 4. But if the statement

Theassumptionthat

2 $\in$Q is false, then its negation is true. This is item 19, which is what we wanted to prove.

We now abandon the robotic, over-detailed proof style adopted in this section, and analyse various methods of proof in more typical settings.

\section{7.2 Proof by Cases}

Some mathematical arguments are made tidier by breaking them into a number of cases, of which precisely one holds, while all lead to the desired conclusion.

Example. Consider the following statement:

The solution set of the inequality 2|x| $\geq$|x $-$1| is the complement of the open interval ($-$1, 1/3).

\begin{proof}
Proof. Let x be a real number. There are three cases.
\end{proof}

(1) If x < 0, then the inequality is $-$2x $\geq$1 $-$x, giving x $\leq$$-$1. (2) If 0 $\leq$x < 1, then the inequality is 2x $\geq$1 $-$x, giving 1/3 $\leq$x < 1. (3) If x $\geq$1, then the inequality is 2x $\geq$x $-$1, giving x $\geq$1.

So the required solution set is the collection of the values of x such that x $\leq$$-$1, or 1/3 $\leq$x < 1, or x $\geq$1, which is the union of two rays:

($-$$\infty$, $-$1] $\cup$[1/3, 1) $\cup$[1, $\infty$) = ($-$$\infty$, $-$1] $\cup$[1/3, $\infty$).

This set is obtained from the real line by removing the open interval ($-$1, 1/3), as claimed. 

The opening sentence `Let x be a real number' acknowledges that all real values of x will have to be considered. The second sentence announces that the proof will branch into three cases, determined by the sign of the expressions within absolute value. Note the careful distinction between strict and non-strict inequalities, to avoid missed or repeated values of x.

The structure of this proof is shaped by the structure of the absolute value function, which is piecewise-defined---see Sect. 5.6. The presence of piecewise-defined functions usually leads to proofs by cases.

Example. A proof by cases on divisibility.

Let n be an integer; then n5 $-$n is divisible by 30.

7.2 Proof by Cases 117

\begin{proof}
Proof. We factor the integer 30 and the polynomial n5 $-$n:
\end{proof}

30 = 2 $\times$ 3 $\times$ 5 n5 $-$n = n(n $-$1)(n + 1)(n2 + 1).

For each prime p = 2, 3, 5 we will show that, for any integer n, at least one factor of n5 $-$n is divisible by p.

1. p = 2. Let n = 2k + j, for some k $\in$Z and j = 0, 1. If j = 0, then n is divisible by 2. If j = 1, then n $-$1 is divisible by 2.

2. p = 3. Let n = 3k + j, for some k $\in$Z and j = 0, $\pm$1. If j = 0, then n is divisible by 3; if j = $\pm$1, then n $\mp$1 is divisible by 3.

3. p = 5. Let n = 5k + j, for some k $\in$Z and j = 0, $\pm$1, $\pm$2. If j = 0, then n is divisible by 5. If j = $\pm$1, then n $\mp$1 is divisible by 5. If j = $\pm$2, then n2 + 1 = 25k2 $\pm$ 20k + 5 = 5(5k2 $\pm$ 4k + 1) is divisible by 5.

The proof is complete. 

After the factorisations we declare our intentions. There are three cases, one for each prime divisor p of 30. Each prime in turn leads to p cases, corresponding to the possible values of the remainder j of division by p. To simplify book-keeping, we consider also negative remainders, and then pair together the remainders which differ by a sign. Note the presence of the symbols $\pm$ and $\mp$within the same sentence. The positive sign in the first expression matches the negative sign in the second expression, and vice-versa.

\section{7.3 Implications}

Many mathematical statements take the form of an implication:

If P then Q.

For example:

If p is an odd prime, then 2p$-$1 $-$1 is divisible by p. If the sequence (ak) is periodic, then (a2k) is also periodic.

Implications may appear in disguise, without the `if\dots{}then' construct:

A is a subset of B. Every repeating decimal is rational. The determinant of an invertible matrix is non-zero.

When we rewrite these sentences as explicit implications, we note that it is necessary to introduce an auxiliary quantity:

118 7 Forms of Argument

If x $\in$A, then x $\in$B. If r is a repeating decimal, then r is rational. If A is an invertible matrix, then det(A) = 0.

This is due to the presence of a hidden universal quantifier. So the symbolic form of an implication is not P $\Rightarrow$Q, but rather

$\forall$x $\in$X, P(x) $\Rightarrow$Q(x) (7.2)

where P and Q are predicates over X. Let us expand one of our sentences so as to tell the full story:

For all real numbers r, if the decimal digits of r are repeating, then r is rational.

Spelling out an implication in this way will help us structure its proof.

\section{7.3.1 Direct Proof}

From the truth table (4.9) we see that if P is false, then P $\Rightarrow$Q is true regardless of the value of Q, so there is nothing to prove. If P is true, then the implication is true provided that Q is true. So a direct proof of the implication consists in assuming

P and deducing Q. The assumption of P lasts only until we have proved Q; after that P is discharged and may no longer be used.

Every direct proof of `If P then Q' must contain a section during which P is assumed. The beginning of the section is announced by a sentence such as

Assume P. Suppose P. Let P.

The task of the rest of the proof is to prove not the original theorem, but Q. We make this clear by writing

RTP: Q

where RTP stands for Remains To Prove (or Required To Prove). The block ends when the proof of Q is completed. This is announced with a closing sentence such as

We have proved Q. The proof of Q is complete.

The above considerations suggest what should be the opening sentence of a direct proof of an implication:

Theorem. If p > 3 is a prime and p+2 is also prime, then p+4 is composite. Proof. Suppose p is a prime number greater than 3, such that p + 2 is prime. RTP: p + 4 is composite.

7.3 Implications 119

All three assumptions (p > 3, p prime, p + 2 prime) must be used in the proof. If not, either the proof is wrong or some assumption is redundant.

Often you can work out how the proof of an implication must start even if you haven't the faintest idea of what the mathematics is about. We illustrate this point with some real life examples:

Theorem. A closed subset of a compact set is compact. Proof. Let X be a compact set, and let C be a subset of X. Assume that C is closed. RTP: C is compact. Theorem. If $\lambda$ $\in$C is a root of a monic polynomial whose coefficients are algebraic integers, then $\lambda$ is an algebraic integer. Proof. Let p be a monic polynomial whose coefficients are algebraic integers, and let $\lambda$ $\in$C be a root of p. RTP: $\lambda$ is an algebraic integer. Theorem. Every basis of a finite-dimensional linear space has the same number of elements. Proof. Let V be a finite-dimensional linear space, and let B1 and B2 be two bases for V . Suppose B1 has n1 elements and B2 has n2 elements. RTP: n1 = n2.

Establishing a good notation is often decisive. In all the examples above, the proof begins by giving names to things. Some authors make this unnecessary by including names in the theorem; others obscure the statement of a theorem by putting too many names in it. Let us rewrite the last theorem in such a way as to establish some notation within the statement itself.

Theorem. Let V be a finite-dimensional linear space. Then every basis for V has the same number of elements. Theorem 1. Let V be a finite-dimensional linear space, and let B1 and B2 be two bases for V . Then \#B1 = \#B2.

Let us compare the three formulations of this theorem. The first is plain and effective. The second contains a minimum of notation, which does no harm but is unnecessary. The last version establishes some useful notation. Normally the notation should be introduced within the proof, unless one plans to use it at a later stage. For instance, one could find the sentence

Let V , B1, B2 be as in Theorem 1.

in the text following the theorem.

\section{7.3.2 Proof by Contrapositive}

In Sect. 4.2 we saw that every implication P $\Rightarrow$Q is equivalent to its contrapositive $\neg$Q $\Rightarrow$$\neg$P: they are both true or both false for any value of P and Q. This equivalence gives us a method for proving implications, called proof by contrapositive, which is a useful alternative to a direct proof. To prove by contrapositive that

120 7 Forms of Argument

If P then Q,

we prove instead that

If not Q then not P,

that is, we assume that Q is false and then deduce that P is false.

A proof by contrapositive is structurally identical to a direct proof, only the predicates are different. Thus in place of (7.2) we consider the statement

$\forall$x $\in$X, $\neg$Q(x) $\Rightarrow$$\neg$P(x). (7.3)

In forming the contrapositive we haven't altered the quantifier; we have merely replaced a boolean function with an equivalent function, much like replacing sin2(x) with 1 $-$cos2(x).

How do we decide between a direct proof and a proof by contrapositive? We must compare the assumptions P and $\neg$Q, and decide which of the two is easier to handle. Sometimes it's necessary to try both approaches to find out. The following examples show the reasoning behind such decisions.

Example. Consider the statement

$\forall$n $\in$N, (2n < n!) $\Rightarrow$(n > 3). (7.4)

The assumption P(n) = (2n < n!) is problematic, because its value is not easily computable. By contrast, $\neg$Q(n) = (n $\leq$3) is straightforward, and the contrapositive implication

$\forall$n $\in$N, n $\leq$3 $\Rightarrow$(2n $\geq$n!) (7.5)

involves checking that the boolean expression (2n $\geq$n!) is true for only three values of n. Proof. We only have to check three cases:

n = 1: 2 = 21 $\geq$1! = 1 n = 2: 4 = 22 $\geq$2! = 2 n = 3: 8 = 23 $\geq$3! = 6.

Thus the expression (7.5) is true, and the proof is complete. 

Example. Consider the statement

If the average of four distinct integers is equal to 10, then one of the integers is greater than 11.

The direct implication involves an assumption on an average value, which entails loss of information; the contrapositive implication involves four integers of bounded size. We opt for the latter, which seems easier.

Given four distinct integers not greater than 11, their average is not equal to 10.

7.3 Implications 121

\begin{proof}
Proof. Let four distinct integers be given. If none of them exceeds 11, then the largest value their sum can assume is 11 + 10 + 9 + 8 = 38. So the largest possible average is
\end{proof}

38 4 = 19

2 < 10

as desired. 

\section{7.4 Conjunctions}

A conjunction is a statement of the type

P and Q.

The statements P and Q are called the conjuncts. As with implications, conjunctions are not necessary presented explicitly.

A direct proof of a conjunction consists of the separate proofs of the two conjuncts, which must be differentiated clearly. It's common practice to put the conjuncts in an ordered list, say (i) and (ii). (If there are more than two conjuncts, the list may be extended using lower-case Roman numerals: (iii), (iv), etc.) Then the proof itself should use the matching labels (i) and (ii), and in each case we should begin by stating what we intend to prove.

\begin{proof}
Proof.
\end{proof}

(i) We prove P. \dots{} (ii) We prove Q. \dots{}

Equality of sets is the canonical hidden conjunction:

A = B.

What are the conjuncts? Two sets are equal if they have the same elements, namely every element of A is an element of B, and vice-versa.

(A = B) $\Leftrightarrow$(A $\subset$B) $\land$(B $\subset$A).

The structure of the proof is now determined by the definition (4.19) of a subset.

\begin{proof}
Proof.
\end{proof}

(i) We prove that A $\subset$B. Let x $\in$A be given. RTP: x $\in$B. (ii) We prove that B $\subset$A. Let x $\in$B be given. RTP: x $\in$A.

Another basic conjunction is the equivalence of two statements:

P if and only if Q.

122 7 Forms of Argument

The equivalence operator $\Leftrightarrow$is the conjunct of an implication and its converse---see (4.12).

\begin{proof}
Proof.
\end{proof}

(i) We prove P $\Rightarrow$Q. \dots{} (ii) We prove Q $\Rightarrow$P. \dots{}

We could replace either part by a proof of the contrapositive implication.

The following theorem is of this type: it provides an alternative characterisation of primality.

Theorem. 3 A natural number n > 1 is prime if and only if n divides (n $-$1)! + 1.

This theorem says P $\land$Q, where

P = $\forall$n > 1, (n is prime) $\Rightarrow$(n | (n $-$1)! + 1) Q = $\forall$n > 1, (n | (n $-$1)! + 1) $\Rightarrow$(n is prime).

The outline of the proof is now clear.

\begin{proof}
Proof.
\end{proof}

(i) Let n be a prime number.

RTP: n divides (n $-$1)! + 1. (ii) Let n be a natural number greater than 1 which divides (n $-$1)! + 1.

RTP: n is prime.

The adverb `precisely' may be used to turn a one-sided implication into an equivalence, hence a conjunction. Thus the proof of the statement

The set Z/nZ is a field precisely if n is prime.

will have to be carried out in two stages.

\begin{proof}
Proof. Let n $\in$N be given.
\end{proof}

(i) Assume that Z/nZ is a field. RTP: n is prime. (ii) Assume that n is prime. RTP: Z/nZ is a field.

\section{7.4.1 Loops of Implications}

The equivalence of several statements is sometimes expressed by a chain of implications in which the last statement in the chain coincides with the first one, thus forming a loop. This type of argument combines conjunctions with implications.

The equivalence of two statements P1 and P2, may be viewed as a loop:

P1 $\Rightarrow$P2 $\Rightarrow$P1.

3 This result, known as Wilson's theorem, was first formulated by Bhaskara (a 7th Century Indian mathematician) and first proved by Lagrange.

7.4 Conjunctions 123

More generally, we consider loops of n implications:

P1 $\Rightarrow$P2 $\Rightarrow$\textperiodcentered{} \textperiodcentered{} \textperiodcentered{} $\Rightarrow$Pn$-$1 $\Rightarrow$Pn $\Rightarrow$P1.

Proving all implications in the loops amounts to proving that all statements are equivalent, namely that Pi $\Leftrightarrow$Pj, for all i, j = 1 . . . , n. This follows from the transitivity of the implication operator.

For example, let G be a group and H a subgroup of G. A left coset of H in G is a set gH = \{gh : h $\in$H\} where g is an element of G. [This is a variant of the algebraic product of sets (2.21)]. A result in group theory states that if x, y are elements of G, then the following statements are equivalent:

(a) x H = yH (b) x H $\subseteq$yH (c) x H $\cap$yH = $\emptyset$ (d) y$-$1x $\in$H.

These are usually proved in a circle, for example (a) $\Rightarrow$(b) $\Rightarrow$(c) $\Rightarrow$(d) $\Rightarrow$(a), but other arrangements are possible, e.g., (a) $\Leftrightarrow$(b) $\Rightarrow$(c) $\Rightarrow$(d) $\Rightarrow$(b).

In Chap. 8 we will use a loop of implications to show the equivalence of four formulations of the principle of induction.

\section{7.5 Proof by Contradiction}

A proof by contradiction of a proposition P consists of assuming $\neg$P and deducing a false statement:

$\neg$P $\Rightarrow$FALSE. (7.6)

The false statement could be anything, including any assertion that contradicts the assumption $\neg$P. Formally, contradiction works because from (7.6) and the truth table (4.9) of the implication operator we deduce that $\neg$P is false and hence P is true.

$\surd$

2, presented in Sect.7.1, we assumed that

For example, in the proof by contradiction of the irrationality of

$\surd$

2 $\in$Q, and deduced that

$\exists$m, n $\in$N, (gcd(m, n) = 1) $\Rightarrow$(m, n $\in$2Z)

which is clearly false.

A classic proof by contradiction is Euclid's proof of the infinitude of primes.

\begin{theorem}
Theorem 7.1. The number of primes is infinite.
\end{theorem}

\begin{proof}
Proof. Assume there are only finitely many primes, p1, . . . , pn. Consider the integer
\end{proof}

n

pk. (7.7)

N = 1 +

k=1

124 7 Forms of Argument

Then N isgreaterthanalltheprimes;moreover,ifwedivide N by pk wegetremainder 1, and therefore N is not divisible by any of the primes. Now, every integer greater than 1 is divisible by some prime. It follows that N must be divisible by a prime not in the given list, a contradiction. 

The statement that every integer greater than 1 is divisible by some prime needs some justification (see Sect. 8.1).

To apply the method of contradiction to an implication P $\Rightarrow$Q, we assume the negation of the implication, namely $\neg$(P $\Rightarrow$Q). By Theorem 4.1, this is equivalent to P $\land$$\neg$Q, so our proof amounts to establishing that

(P $\land$$\neg$Q) $\Rightarrow$FALSE.

This form of proof by contradiction is called the both ends method: to prove that if P then Q, we assume both P and not-Q, and we deduce something impossible. The appealing feature of this method is that it gives us two assumptions to exploit, namely P and $\neg$Q. This is often how mathematicians work, although this strategy may not be made explicit. It also easily causes confusion in writing, as the following example (from a textbook) shows.

Theorem. For all real numbers x, if x > 0, then 1/x > 0.

Bad Proof: If 1/x < 0 then ($-$1)/x > 0, so x(($-$1)/x) > 0, i.e., $-$1 > 0, a contradiction.Therefore1/x $\geq$0.Since1/x can'tbe0,weconcludethat1/x > 0. 

Inthisproof,theassumptionof P ishidden,andtheassumptionof$\neg$Q isconfused. Not helpful writing!

\begin{proof}
Proof. We use contradiction. Let us assume that x > 0 and 1/x $\leq$0. Multiplying the second inequality by $-$1, we obtain $-$1/x $\geq$0, and since x is positive, multiplication by x yields x($-$1/x) $\geq$0. Simplification gives $-$1 $\geq$0, which is false. 
\end{proof}

The final deduction rests on the inequality 0 < 1, which may be derived from the axioms of the real number system.

\section{7.6 Counterexamples and Conjectures}

Suppose that a statement concerning all elements of a set is false. To prove this, it is sufficient to exhibit a single element of that set for which the statement fails. This argument is called a counterexample. In symbols, we show that the statement

$\forall$x $\in$X, P(x)

is false by proving its negation, namely

$\exists$x $\in$X, $\neg$P(x).

7.6 Counterexamples and Conjectures 125

To explore the idea of a counterexample, let us consider a famous polynomial proposed by Euler: p(n) = n2 + n + 41. We evaluate p(n) at integer values of the

indeterminate: n = 0, 1, \dots{}

n 0 1 2 3 4 5 6 7 8 9 10 \textperiodcentered{} \textperiodcentered{} \textperiodcentered{} 20 \textperiodcentered{} \textperiodcentered{} \textperiodcentered{} 30 p(n) 41 43 47 53 61 71 83 97 113 131 151 \textperiodcentered{} \textperiodcentered{} \textperiodcentered{} 461 \textperiodcentered{} \textperiodcentered{} \textperiodcentered{} 971

All these values of p(n) are prime! Moreover, p($-$n) = p(n $-$1), so it seems plausible---if a bit daring---to put forward the following conjecture:

Conjecture 1. For all integers n, p(n) is prime.

A conjecture is a statement that we wish were a theorem. Three things may happen to a conjecture: (i) someone proves it, and the conjecture becomes a theorem; (ii) someone produces a counterexample, and the conjecture is proved false; (iii) none of the above, and the conjecture remains a conjecture.

Our case is (ii). The negation of conjecture 1 is

There is an integer n such that p(n) is composite.

To prove it, we exhibit such a value of n. For n = 40, we find

p(40) = 402 + 40 + 41 = 40(40 + 1) + 41 = 41(40 + 1) = 412.

So p(40) is not prime, and conjecture 1 is false.

For a mathematician, the formulation of a conjecture is accompanied by the fear of a counterexample. Let us return to conjecture 1: we know it's false, but can it be modified into a weaker statement? It can be verified that n = 40 is the smallest value for n for which p(n) is not prime.4 Are there other such values? If these values were exceptional, then a meaningful weaker version of conjecture 1 could be

Conjecture 2. For all but finitely many integers n, p(n) is prime.

Even if p(n) were composite for, say, a billion values of n, the conjecture would still hold. Unfortunately, conjecture 2 is not true either. Its negation

There are infinitely many integers n for which p(n) is composite.

is established by the simple, devastating counterexample:

p(41k) = (41k)2 + 41k + 41 = 41(41k2 + k + 1) k = 1, 2, . . . (7.8)

which says that 41 is a proper divisor of p(41k) for infinitely many values of k.

Let's make one final attempt to salvage something from our original claim.

Conjecture 3. There are infinitely many integers n such that p(n) is prime.

4 This phenomenon has deep roots, see [9, p. 155].

126 7 Forms of Argument

This is a very meaningful conjecture. Nobody has ever been able to prove that any quadratic polynomial with integer coefficients assumes infinitely many prime values. On the other hand, if we perform a numerical experiment we discover that, as n increases, p(n) continues to supply prime values, even though the composite values slowly become dominant (see Exercise 10.4.1). So this conjecture is very likely to be true, and it makes sense to put it forward.

Arguably, the most famous conjecture in mathematics is the Riemann's hypothesis (RH), formulated in 1859. This conjecture (essentially) says that the zeta-function defined in Eq. (6.3), vanishes only at points whose real part is equal to 1/2. The depth of this statement is not at all apparent from its idiosyncratic formulation.

In a sense, the mathematics community cannot afford to wait for this important matter to be settled, and several theorems about the RH have been proved. Some of these theorems formulate the RH in terms of equivalent statements, establishing the importance of the RH in many areas of mathematics. There are also proofs of conditional theorems, which assume the validity of the RH, and deduce some of its consequences.

Other famous conjectures include the twin-primes conjecture, attributed to Euclid:

There are infinitely many primes p such that p + 2 is also prime. and the Goldbach conjecture5:

Every even integer greater than 2 can be written as a sum of two primes. Both conjectures are easy to formulate and widely believed to be true, but their proof seems hopelessly difficult. This situation is common in the theory of numbers, where assessing the difficulty of a problem is not always easy. To illustrate this phenomenon, let us arrange the natural numbers into four columns, with the primes highlighted in boldface:

1 2 3 4 5 6 7 8 9 10 11 12 13 14 15 16 17 18 19 20 21 22 23 24 25 26 27 28 29 39 31 32 33 34 35 36 37 38 39 40 41 42 43 44 45 46 47 48 49 50 51 52 53 54 55 56 57 58 59 60

\dots{}

5 Christian Goldbach (German: 1690--1764).

7.6 Counterexamples and Conjectures 127

Weguidethereaderthroughthesedatabyaskingsomequestions,arrangedinorder of increasing difficulty. (We have considered this form of exposition in Sect. 6.5.)

1. Why is there just one prime in the second column and no primes at all in the

fourth?

2. The first two rows combined contain four primes; are there other pairs of adjacent

rows with the same property?

3. Does the table contain infinitely many primes?

4. There are rows without primes. Are there infinitely many of them?

5. Are there adjacent rows without primes? What about sequences of three or more

adjacent rows without primes?

6. Are there infinitely many primes in the first column?

7. Are `half' of the primes in column 1 and `half' in column 3?

8. Are there infinitely many rows containing two primes?

At the end of primary school, a pupil should be able to answer question 1, while in secondary school one could see why the answer to question 2 must be negative (one of three consecutive odd integers is divisible by 3). The answer to questions 3 to 7 is affirmative. Euclid's theorem (Sect. 7.5), which answers question 3, is normally taught in a first year university course, although we have seen that the proof does not require advanced ideas. Answering questions 4 and 5 does involve some ingenuity, but no major difficulty.

The mathematics needed to answer 6 and 7 becomes difficult. The proof that column 1 contains infinitely many primes, known since the 17th Century, is given today in an undergraduate course in number theory. The formulation of question 7 can be made precise, and the affirmative answer was given at the end of the 19th Century. Understanding the proof requires a solid background in analysis.

Many mathematicians would conjecture that the answer to question 8 is again positive, but at present nobody knows how to prove it. This is a variant of the twin primes conjecture stated above.

An accessible introduction to this beautiful part of mathematics is found in [38].

\section{7.7 Wrong Arguments}

In constructing a mathematical argument it's easy to make mistakes. In this section we identify some common faulty arguments: confusing examples with proofs, assuming what we are supposed to prove, mishandling functions. Other mistakes will be examined in Chap. 9 and in the exercises. Awareness of these problems should help us in avoiding them.

\section{7.7.1 Examples Versus Proofs}

The verification of a statement in specific cases does not constitute a form of proof. Our study of Euler's polynomial in Sect. 7.6 shows how misleading examples can be.

128 7 Forms of Argument

This state of affairs is peculiar to mathematics; in other scientific disciplines, such as biology, a proof of a statement consists of independent experimental verifications of its validity.

In our first example the fault is easy to spot.

Theorem. For all primes p, the integer 2p $-$2 is divisible by p.

Wrong Proof.

22 $-$2 = 2 \textperiodcentered{} 1, 23 $-$2 = 3 \textperiodcentered{} 2, 25 $-$2 = 5 \textperiodcentered{} 6, 27 $-$2 = 7 \textperiodcentered{} 18, etc. 

The theorem has been proved only for p = 2, 3, 5, 7.

The next example is similar, but not so clear-cut [36, p. 138f].

Theorem. For all x, y and z, if x + y < x + z then y < z.

Wrong Proof. Suppose x + y < x + z. Take x = 0. Then

y = 0 + y < 0 + z = z. 

The mistake here is that we took x to be 0, which is a special value of x. This assumption is wholly unjustified, since the quantities x, y, z are controlled by an existential quantifier, and hence no condition may be imposed on them. By adding the assumption that x = 0, we have in fact proved the

Weaker Theorem. For all y and z, if 0 + y < 0 + z, then y < z.

This is not what we claimed to prove.

\section{7.7.2 Wrong Implications}

Inappropriate handling of implications results in common mistakes. Instead of proving an implication, we may end up proving its converse; or we may assume the statement we are meant to prove, deduce from it a true statement, and believe we've completed the proof. These faulty deductions---of which we now show an example--- are sometimes called circular arguments, or non sequiturs.6

$\surd$ 6 <

$\surd$

$\surd$ 2 +

Prove that

15

Wrong Proof.

$\surd$ 6 <

$\surd$ 2 +

$\surd$ 15 $\Rightarrow$(

$\surd$ 2 +

$\surd$

6)2 < 15

$\surd$

$\Rightarrow$8 + 2

12 < 15

$\surd$

$\Rightarrow$2

12 < 7

$\Rightarrow$48 < 49. 

6 Latin for `it does not follow'.

7.7 Wrong Arguments 129

$\surd$ 6 <

$\surd$

$\surd$ 2 +

15). Instead we have assumed P, and correctly deduced from it the true statement 48 < 49. However, the deduction P $\Rightarrow$TRUE (unlike the deduction P $\Rightarrow$FALSE) gives us no information about P, from Table (4.9). Indeed, had we started from the false statement

We were supposed to prove P, where P = (

$\surd$

$\surd$ 2+

6 < $-$

$\surd$

15, we would have reached exactly the same conclusion. There are two methods for fixing this problem.

First method: retracing the steps. We regard the chain of deductions displayed above as `rough work'; then we start from the end and prove the chain of converse implications.

\begin{proof}
Proof.
\end{proof}

$\surd$ 48 < 49 $\Rightarrow$

$\surd$ 48 <

49

$\surd$

$\Rightarrow$2

12 < 7 $\Rightarrow$8 + 2

$\surd$

12 < 15

$\surd$

$\surd$ 2 +

6)2 < 15 $\Rightarrow$

$\Rightarrow$(

$\surd$ 6 <

$\surd$

$\surd$ 2 +

15

where in the first and the last implications we have taken the positive square root of each side. We have proved the implication TRUE $\Rightarrow$P, from which we deduce that

P is TRUE. 

Clearly, we could not have come up with such a proof without having done the `rough work' first; the next method does not require this.

Second method: contradiction.

\begin{proof}
Proof. Let us assume $\neg$P:
\end{proof}

$\surd$ 6 $\geq$

$\surd$

$\surd$ 2 +

$\surd$ 15 $\Rightarrow$(

$\surd$ 2 +

6)2 $\geq$15

$\surd$

$\Rightarrow$8 + 2

12 $\geq$15 $\Rightarrow$2

$\surd$

12 $\geq$7

$\Rightarrow$48 $\geq$49.

We have proved that $\neg$P $\Rightarrow$FALSE, which implies that $\neg$P is FALSE, that is, P is TRUE. 

\section{7.7.3 Mishandling Functions}

Many things can go wrong when we deal with functions. For instance, we may apply a function to arguments outside its domain (e.g., forming the logarithm of zero), or invert a function incorrectly (e.g., taking the wrong sign of a square root).

We begin with a proof from calculus which isn't correct, even if it captures the essence of the argument. The chain rule for differentiating the composition of two

130 7 Forms of Argument

functions states that if f and g are differentiable, then

(g f )$'$(x) = g$'$( f (x)) f $'$(x).

Sloppy Proof. Let $\varepsilon$ be a real number. We compute

g( f (x + $\varepsilon$)) $-$g( f (x))

= g( f (x + $\varepsilon$)) $-$g( f (x)) $\varepsilon$

$\times$ f (x + $\varepsilon$) $-$f (x) f (x + $\varepsilon$) $-$f (x)

. $\varepsilon$

Letting $\varepsilon$ tend to zero gives the desired result. 

The right-hand side is obtained by multiplying and dividing by f (x + $\varepsilon$) $-$f (x). But this quantity may well be zero for non-zero $\varepsilon$ (e.g., if f is constant), which invalidates the argument. A valid proof requires more delicacy.

\begin{proof}
Proof. Let y = f (x). Since g is differentiable at y, we write
\end{proof}

g(y + $\delta$) $-$g(y) = g$'$(y)$\delta$ + h(y, $\delta$)$\delta$

where h is a continuous function with h(y, 0) = 0. Specialising to $\delta$ = f (x + $\varepsilon$) $-$

f (x), we find

= g$'$( f (x)) f (x + $\varepsilon$) $-$f (x) $\varepsilon$

g( f (x + $\varepsilon$)) $-$g( f (x))

$\varepsilon$

+ h(y, f (x + $\varepsilon$) $-$f (x)) f (x + $\varepsilon$) $-$f (x)

. $\varepsilon$

As $\varepsilon$ tend to zero, the last term tends to zero, being the product of a function that tends to zero and a function that tends to a finite limit (because f is differentiable). This gives the desired result. 

Mishandling a non-invertible function may have unpredictable consequences:

Wrong Theorem. Every negative real number is positive.

Wrong Proof. Let $\alpha$ be a negative real number. Then

$\alpha$ = $-$1|$\alpha$| = ($-$1)1|$\alpha$| = ($-$1)2\textperiodcentered{} 1

2 |$\alpha$| = [($-$1)2] 2 |$\alpha$| 1

1 2 |$\alpha$| = |$\alpha$| > 0

= 1

as required. 

What's wrong with this proof? The function x $\to$x2 is not invertible, and to the positive sign for the square root, then the functions x $\to$x2 and x $\to$$\surd$x are invert it we must restrict the domain. If we restrict it to the range x $\geq$0 and take inverse of each other: x2\textperiodcentered{} 1

$\surd$ 2 =

x2 = x. However, if x is negative, these functions are inverse of each other only if we take the negative sign of the square root.

Thus the error occurred in the last equality

7.7 Wrong Arguments 131

1 2 = 1 1

where, having squared a negative quantity, we mistakenly choose the positive sign of the square root. Exercise 7.3.2 deals with a more general instance of the same phenomenon.

\section{7.8 Writing a Good Proof}

Proofs come in all shapes and sizes. Correctness is, of course, imperative, but a good proof requires a lot more since the conflicting requirements of clarity and conciseness must be resolved. On the one hand, reading a proof is demanding, and any assistance will be welcome. On the other hand, a proof is a rigorous exercise in logic and the exposition must remain succinct and essential. The level of exposition and the amount of detail will depend on the mathematical maturity of the target audience.

There is one general rule, to be applied at all key junctures in a proof:

\textbullet{} Say what you plan to do.

\textbullet{} When you've done it, say so.

We now examine statements of theorems and proofs that are less than optimal, and we seek to improve them. Our proofs are very detailed, and suitable for first-year university students. In some cases we also provide a concise version of the proof, aimed at expert readers.

Our first theorem is taken from a textbook.

Bad Theorem. If x2 = 0, then x2 > 0.

Bad proof. If x > 0 then x2 = xx > 0. If x < 0 then $-$x > 0 , so ($-$x)($-$x) > 0, i.e., x2 > 0. 

The statement of the theorem is incomplete, in that there is no information about the quantity x. The theorem is an implication, but in the proof what has happened to the assumption x2 = 0? If the assumption is obvious, one may leave it out, but one must judge whether the readers are sophisticated enough to see what it must be. Furthermore, this is a proof by cases (see Sect. 7.2), and it would be helpful if this were made explicit.

Theorem. For all real numbers x, if x2 = 0, then x2 > 0.

\begin{proof}
Proof. Let x be a real number such that x2 = 0. Then x = 0, and we have two cases:
\end{proof}

(i) x < 0. Then $-$x > 0 , so ($-$x)($-$x) > 0, that is, x2 > 0. (ii) x > 0. Then xx = x2 > 0. 

The following statement establishes a simple relation between rational and irrational numbers.

132 7 Forms of Argument

Bad Theorem. $\forall$a, b $\in$R, (a $\in$Q $\land$b $\in$Q) $\Rightarrow$a + b $\in$Q.

Bad Proof: Suppose a + b $\in$Q. Then a + b = m/n. If a $\in$Q, then a = p/q, and b = m/n $-$p/q $\in$Q. 

The symbolic formulation of the theorem obscures its simple content. The proof is straightforward, but the use of unnecessary symbols (the numerator and denominator of the rational numbers serve no purpose) and the lack of comments make it less clear than it should be.

Theorem. The sum of a rational and an irrational number is irrational.

\begin{proof}
Proof. Let a and z be a rational and an irrational number, respectively. Consider the identity z = (a + z) $-$a. If a + z were rational, then z, being the difference of two rational numbers, would also be rational, contrary to our assumption. Thus a + z must be irrational. 
\end{proof}

Thechoiceofsymbolskeepstherationalandtheirrationalnumberswell-separated in our mind.

Next we consider an arithmetical statement.

Bad Theorem. n odd $\Rightarrow$8|n2 $-$1.

Bad proof. n odd $\Rightarrow$$\exists$j $\in$Z, n = 2 j + 1; $\therefore$n2 $-$1 = 4 j( j + 1); $\forall$j $\in$Z, 2 | j( j + 1) $\Rightarrow$8 | n2 $-$1 

This is a clumsy attempt to achieve conciseness via an entirely symbolic exposition.Combiningwordsandsymbolsandaddingsomeshortexplanationswillimprove readability and style. We shall see that a clearer proof need not be longer.

Theorem. The square of an odd integer is of the form 8n + 1.

Detailed Proof. Let k be an odd integer. We show that k2 = 8n + 1, for some n $\in$Z.

Since k is odd, we have k = 2 j + 1, for some integer j. We find

k2 = (2 j + 1)2 = 4 j2 + 4 j + 1 = 4 j( j + 1) + 1. (1)

Now, one of j or j + 1 is even, and therefore their product is even. Thus we have

j( j + 1) = 2n, for some n. Inserting this expression in (1) we obtain k2 = 8n + 1, as desired. 

This proof is quite detailed, and can safely be shortened.

\begin{proof}
Proof. An odd integer k is of the form k = 2 j + 1, for some integer j. A straight-forward manipulation gives k2 = 4 j( j + 1) + 1. Our claim now follows from the observation that the product j( j + 1) is necessarily even. 
\end{proof}

The adverb `necessarily' signals that the conclusion (the product j( j +1) is even) requires a moment's thought. The expression `straightforward manipulation' is used

7.8 Writing a Good Proof 133

to omit some steps in the derivation, while warning the reader that some calculations are required. In such a circumstance, the adjective `straightforward' is preferable to `easy', which is subjective, or `trivial', which carries a hint of arrogance. So it's appropriate to claim that the arithmetical identity

$\surd$

$\surd$

$\surd$$-$1 (1 + 2 +

$\surd$$-$1 (1 + 2 $-$

5) 2

5) 2

29 =

$\surd$

$\surd$

$\surd$$-$1 (1 $-$ 2 $-$

$\surd$$-$1 (1 $-$ 2 +

5) 2

5) 2

$\times$

can be verified with a `straightforward calculation'.

In the following example everything needs re-writing: definition, theorem, and proof.

Bad Definition. Let q j be digits. We define

q1 \textperiodcentered{} \textperiodcentered{} \textperiodcentered{} qn = q1 \textperiodcentered{} \textperiodcentered{} \textperiodcentered{} qnq1 \textperiodcentered{} \textperiodcentered{} \textperiodcentered{}

Bad Theorem. Let x = m + 0.q$'$

1 \textperiodcentered{} \textperiodcentered{} \textperiodcentered{} q$'$ kq1 \textperiodcentered{} \textperiodcentered{} \textperiodcentered{} qn, where m $\in$Z. Then x $\in$Q.

Bad Proof: Let A = 0.q$'$

1 \textperiodcentered{} \textperiodcentered{} \textperiodcentered{} q$'$ k, B = 0.q1 \textperiodcentered{} \textperiodcentered{} \textperiodcentered{} qn. Then A, B $\in$Q, and

10n = B

1 1 1 + 10n + 102n + \textperiodcentered{} \textperiodcentered{} \textperiodcentered{}

0.q1 . . . qn = B

10n $-$1 =: C $\in$Q.

Thus x = m + A + C10$-$k $\in$Q. 

In the definition, the decimal nature of the digits is not mentioned and the periodicity of the digit sequence could be made more evident. In the statement of the theorem the number x is represented with an inappropriate hybrid notation, as the sum of a decimal number and a number expressed symbolically. The proof lacks explanations and the choice of symbols is not ideal.

As we did before, we state the theorem with words, which is more effective. We introduce the notation after the statement of the theorem but before the beginning of the proof. This arrangement should be considered if the notation is needed outside the confines of the proof, or if the proof is heavy and needs lightening up, or simply as a variation from the rigid definition-theorem-proof format. We number the various steps in the argument for future reference.

Theorem. Every number whose decimal digits eventually repeat is rational.

1. Before proving the theorem, we establish some notation. We denote a string (finite orinfinite)ofdecimaldigitsbyd1d2 \textperiodcentered{} \textperiodcentered{} \textperiodcentered{} ,withdk $\in$\{0, 1, . . . , 9\};theover-bardenotes a pattern of digits which repeats indefinitely:

d1 \textperiodcentered{} \textperiodcentered{} \textperiodcentered{} dn = d1 \textperiodcentered{} \textperiodcentered{} \textperiodcentered{} dnd1 \textperiodcentered{} \textperiodcentered{} \textperiodcentered{} dn \textperiodcentered{} \textperiodcentered{} \textperiodcentered{} . (7.9)

134 7 Forms of Argument

Detailed Proof.

2. Let x be a real number with eventually repeating decimal digits. Then x has a decimal representation of the type

x = d$'$

0 \textperiodcentered{} \textperiodcentered{} \textperiodcentered{} d$'$ j.d$'$

j+1 \textperiodcentered{} \textperiodcentered{} \textperiodcentered{} d$'$

md1 \textperiodcentered{} \textperiodcentered{} \textperiodcentered{} dn

for some m $\geq$0, 0 $\leq$j $\leq$m, and n $\geq$1. The primed digits form the aperiodic part of the digit sequence.

3. We shall compute the value of x in terms of two integers M and N, given by:

m

n

M = d$'$

0 \textperiodcentered{} \textperiodcentered{} \textperiodcentered{} d$'$ m =

d$'$

k10m$-$k N = d1 \textperiodcentered{} \textperiodcentered{} \textperiodcentered{} dn =

dk10n$-$k.

k=0

k=1

First, we get rid of the aperiodic part; we shift it to the left of the decimal point by multiplying by a suitable power of 10, and then we subtract it off, to get

10m$-$j x $-$M = 0.d1 \textperiodcentered{} \textperiodcentered{} \textperiodcentered{} dn. (7.10)

4. Next we compute, using (7.9)

0.d1 \textperiodcentered{} \textperiodcentered{} \textperiodcentered{} dn = N N N

10n + 102n + 103n + \textperiodcentered{} \textperiodcentered{} \textperiodcentered{}

1

1 1 10n + 102n + 103n + \textperiodcentered{} \textperiodcentered{} \textperiodcentered{}

= N

1

k

$\infty$

N = 10n $-$1.

= N

10n

k=1

5. In the last step we have used the formula of the sum of the geometric series

$\infty$

qk = q |q| < 1.

1 $-$q

k=1

From Eq. (7.10) and the result above, we obtain

10m$-$j x $-$M = N

10n $-$1

and hence

N M + 10n $-$1

x = 10 j$-$m

.

The right-hand side of this equation consists of sums and products of rational numbers, and is therefore rational. 

7.8 Writing a Good Proof 135

Let us examine the main steps in the proof:

1. Since we don't begin the proof straight away, we say so. The key notation (the

over-bar) is established first in words, then in symbols. The periodicity of the

digit sequence is clearer.

2. In the theorem there is a hidden quantifier, and the proof begins accordingly. Then

we introduce the notation for x, and we say why.

3. At every opportunity, we declare our intentions. A coherent choice of symbols

improves readability.

4. In this passage, one must decide what constitutes an appropriate amount of details:

we have been cautious here.

5. A reminder of a well-known summation formula.

We give a concise version of the same proof.

\begin{proof}
Proof. Let x be a real number with eventually repeating decimal digits dk. Without loss of generality, we may assume that the integer part of x is zero, and that the fractional part is purely periodic:
\end{proof}

x = 0.d1 \textperiodcentered{} \textperiodcentered{} \textperiodcentered{} dn. (1)

(Any real number may be reduced to this form by first multiplying by a power of 10, and then subtracting an integer, and neither operation affects the tail of the digit sequence.) Defining the decimal integer M = d1 \textperiodcentered{} \textperiodcentered{} \textperiodcentered{} dn, we find, from (1)

1

k

$\infty$

x = M

M M 10n + 102n + 103n + \textperiodcentered{} \textperiodcentered{} \textperiodcentered{} = M

M = 10n $-$1. (7.11)

10n

k=1

We see that x is rational. 

The expression `without loss of generality' indicates that the restriction being introduced does not weaken the argument in any way (see also Sect. 5.2). A brief parenthetical remark is added for clarification: it too could be omitted.

Exercise 7.1 You are given cryptic proofs of mathematical statements. Rewrite them in a good style, with plenty of explanations.

1. Provethat thelinethroughthepoint (4, 5, 1) $\in$R3 parallel tothevector (1, 1, 1)T

and the line through the point (5, $-$4, 0) parallel to the vector (2, $-$3, 1)T intersect at the point (1, 2, $-$2). (The symbol T denotes transposition.) Bad proof.

(4, 5, 1)T + $\lambda$ (1, 1, 1)T = (4 + $\lambda$, 5 + $\lambda$, 1 + $\lambda$)T ;

= (5, $-$4, 0)T + $\mu$ (2, $-$3, 1)T = (5 + 2$\mu$, $-$4 $-$3$\mu$, $\mu$)T . $\therefore$4 + $\lambda$ = 5 + 2$\mu$, 5 + $\lambda$ = $-$4 $-$3$\mu$, 1 + $\lambda$ = $\mu$. $\Rightarrow$4 + $\lambda$ = 5 + 2(1 + $\lambda$) $\Rightarrow$$\lambda$ = $-$3, $\mu$ = $-$2. $\therefore$v = (1, 2, $-$2)T . 

136 7 Forms of Argument

2. Prove that the real function x $\to$3x4 + 4x3 + 6x2 + 1 is positive.

Bad proof.

f $'$(x) = 12x(x2 + x + 1) = 0 $\Leftrightarrow$x = 0; f $'$$'$(0) > 0, f (0) > 0 $\Rightarrow$positive. 

3. Prove that for all real values of a, the line

a $-$1

2

y = ax $-$

2

is tangent to the parabola y = x2 + x. Bad proof.

ax $-$(a $-$1)2/4 = x2 + x; $\Rightarrow$x2 + x(1 $-$a) + (a $-$1)2/4 = 0; $\Rightarrow$x = (a $-$1)/2. For this x, dy/dx is the same. 

4. Prove that the vector (1, 0, 1) does not belong to the subspace of R3 generated by the vectors (3, 1, 2) and ($-$2, 1, 0). Bad proof:

3t $-$2s = 1; t + s = 0; 2t = 1. $\Rightarrow$t = 1/2 $\Rightarrow$s = $-$1/2 $\Rightarrow$3/2 = 0, contradiction. 

Exercise 7.2 The following text has several faults: (a) explain what they are; (b) write an appropriate revision. Wrong theorem. For all numbers x and y,

x2 + y2

> 2. |xy|

Wrong proof. For

x2 + y2

> 2 $\Rightarrow$x2 + y2 > 2xy |xy|

$\Rightarrow$x2 $-$2xy + y2 > 0 $\Rightarrow$(x $-$y)2 > 0.

The last equation is trivially true, which proves it. 

Exercise 7.3 Identify the flaw in the proof of the following statements.

1. The integral of a positive function may be negative.

Wrong proof. The function f (x) = 1/x2 is positive. We compute

b

b

dx

a $-$1

$-$1

= 1

b.

x2 =

I (a, b) =

x

a

a

7.8 Writing a Good Proof 137

We find that I ($-$1, 1) = $-$2 < 0. 

2. The complex exponential function is constant. Wrong proof. Let $\theta$ be an arbitrary complex number. We find

e2 $\pi$ i $\theta$

2 $\pi$ = 1 2 $\pi$ = 1. $\theta$ 

ei$\theta$ = ei$\theta$ 2 $\pi$

2 $\pi$ = e2 $\pi$ i $\theta$ 2 $\pi$ =

Exercise 7.4 Write the first few sentences of the proof of each statement, introducing all relevant notation in an appropriate order, and identifying the RTP. (For this task, understanding the meaning of the statements is unimportant.)

1. On a compact set every continuous function is uniformly continuous.

2. Every valuation of a field with prime characteristic is non-archimedean.

3. No polynomial with integer coefficients, not a constant, can be prime for all integer values of the indeterminate.

4. Any hyperbolic matrix in SL2(Z) is conjugate to a matrix with non-negative coefficients.

5. The only bi-unique analytic mapping of the interior of the unit disc into itself is a linear fractional transformation.

6. In a field with a non-archimedean absolute value, the connected component of any point is the set consisting of only that point.

7. The characteristic polynomial of the incidence matrix of a Pisot substitution is irreducible over Q.

8. The first Betti number of a compact and orientable Riemannian manifold of positive Ricci curvature is zero.

9. Every differential of the second or third kind differs from some normal differential by a differential of the first kind.

10. A real function continuous in a closed interval attains all values between its maximum and minimum.

11. A subset of a metric space is open if and only if its complement is closed.

12. A stable set of a graph is maximum if and only if there exists no maximal alternating sequence of odd length.

13. A bi-infinite sequence over the alphabet \{0, 1\} is Sturmian if and only if it is balanced and not eventually periodic.

14. The curvature of a hermitian manifold vanishes if and only if it is possible to choose a parallel field of orthonormal holomorphic frames in a neighbourhood of each point of the manifold.

Exercise 7.5 You are asked to help the reader approach the following problem:

Prove that for all integers m, n, we have mZ + nZ = gcd(m, n)Z.

Formulate a number of easier sub-problems of increasing difficulty, eventually leading to the original problem.

\end{document}